% Complete proofs for the fixed algorithms; historical source obligations remain separate.
\begin{annexitem}{annex:p07.combinatory.bracket-abstraction}{thm:p07-bracket-abstraction}{The original host bracket-abstraction property}
\begin{proof}
The host algorithm of Definition~\ref{def:p07-sk-combinators} is exactly
\(\mathcal B\) in Definition~\ref{def:p07-combinatory-abstraction-algorithms}.
Theorem~\ref{thm:p07-combinatory-substitution}, with \(Q=x\), proves that
\(\mathcal B_x(P)\) contains no \(x\) and
\(\mathcal B_x(P)x\longrightarrow_{\mathsf{SK}}^*P[x/x]=P\).
Recursive binder elimination is the finite map \(T_{\mathcal B}\) of
Definition~\ref{def:p07-combinatory-translations}.
Theorem~\ref{thm:p07-combinatory-translation} proves its precise correspondence:
beta-correct lambda expansion and weak simulation of root beta contractions
in application contexts.  Theorem~\ref{thm:p07-combinatory-extensional} gives
the full conversion correspondence when extensional equality is specified.
These facts establish the bracket-abstraction property and the stated
translation on the fixed S--K basis.
\end{proof}
\end{annexitem}

\begin{annexitem}{annex:p07.combinatory.substitution}{thm:p07-combinatory-substitution}{All abstraction cases and substitution side conditions}
\begin{proof}
Every recursive call uses a proper subtree, so induction proves termination.
At input \(x\), the output I has no variables and
\(\mathsf I Q=\mathsf S\mathsf K\mathsf KQ\) contracts to
\(\mathsf KQ(\mathsf KQ)\), then to \(Q\).
For \(\mathcal B\), an input \(P\) without \(x\) gives \(\mathsf KP\),
with exactly the variables of \(P\), and \(\mathsf KPQ\) contracts to
\(P=P[Q/x]\).  For \(\mathcal A\), the other-variable and constant cases
give \(\mathsf KyQ\longrightarrow_{\mathsf{SK}}y\) and
\(\mathsf KcQ\longrightarrow_{\mathsf{SK}}c\), respectively.
For an application \(UV\), the induction hypotheses give
\[
 \mathsf S\mathcal D_x(U)\mathcal D_x(V)Q
 \longrightarrow_{\mathsf{SK}}
 (\mathcal D_x(U)Q)(\mathcal D_x(V)Q)
 \longrightarrow_{\mathsf{SK}}^*U[Q/x]V[Q/x].
\]
The variable set is the union of the recursive variable sets, namely
\(\FV(UV)\setminus\{x\}\).  This proves both assertions in every case.

Assume \(y\ne x\) and \(x\notin\FV(Q)\).  For host commutation with
\([Q/y]\), input \(x\) gives I on both sides.  Substitution preserves absence
of \(x\), so the absent case gives \(\mathsf K(P[Q/y])\) on both sides.
If an application contains \(x\), substitution for \(y\) neither removes nor
introduces \(x\); both sides use the application clause, and induction on its
children proves literal equality.
For \(\mathcal A\), input \(x\) again gives I on both sides.
Constants and variables different from \(x,y\) are
unchanged.  At \(y\), substitution gives \(\mathsf KQ\), whereas recomputation
gives \(\mathcal A_x(Q)\).  For an atom \(Q\) these are the same tree.
For arbitrary \(Q\) independent of \(x\), both terms applied to a fresh variable
reduce to \(Q\), so they are extensionally equal.  Application congruence and
induction prove the stated extensional commutation.  This argument uses only
the application theorem and the defining fresh-variable equality rule.

For fresh renaming \([z/x]\), the identity case remains the identity case,
each \(x\)-free subtree stays unchanged, constants and other atoms stay unchanged,
and induction applies to the children of an application.  Hence both algorithms
satisfy the literal renaming identity.
Every S--K contraction is stable under arbitrary tree substitution: verify
this directly in \(\mathsf KUV\longrightarrow U\) and
\(\mathsf SUVW\longrightarrow UW(VW)\), then pass through application contexts.
Apply the one-variable theorem to the outermost abstractor in the iterated
construction.  Apply it next to the following abstractor, substituting the
arguments already supplied into that next reduction sequence.  This is a
reduction of the constructed tree even when recomputing an abstractor would
give a different tree.  The freshness of the \(Q_i\) makes each successive
replacement leave every previously inserted argument unchanged.  Induction on
\(n\) gives exactly the displayed simultaneous substitution; \(n=0\) is identity.
\end{proof}
\end{annexitem}

\begin{annexitem}{annex:p07.combinatory.translation}{thm:p07-combinatory-translation}{Capture avoidance, expansion, and application simulation}
\begin{proof}
Expansion preserves free variables and commutes with substitution by induction,
choosing the bound variables in its closed constant expansions fresh.
For \(E(\mathcal D_x(P))=_\beta\lambda x.E(P)\), induct on \(P\).
The expansion of I reduces to \(\lambda x.x\).  Each constant-output case
\(E(\mathsf KP)\) reduces to \(\lambda x.E(P)\), with \(x\) fresh for
\(E(P)\).  This includes the other-variable and constant cases of \(\mathcal A\).
For an application put \(U'=E(\mathcal D_x(U))\) and
\(V'=E(\mathcal D_x(V))\).  Its expansion reduces to
\(\lambda z.U'z(V'z)\), for fresh \(z\).  The induction hypotheses make this
beta-equal to \(\lambda z.E(U)[z/x]E(V)[z/x]\), which is alpha-equal to
\(\lambda x.E(UV)\).

The compilers preserve free variables by induction using the abstraction
variable equality.  Compilation commutes with substitution of a variable:
variables and application are immediate, and at a lambda binder first choose
its name fresh and use atomic abstraction commutation.  A fresh renaming of a
binder therefore renames the compiled body correspondingly; the abstraction
renaming identity makes the output literally the same.  Induction through
outer constructors and finite sequences of renamings proves alpha invariance.

For arbitrary substitution under \(T_{\mathcal B}\), induct on \(M\).
Variables and application are immediate.  At \(\lambda y.R\), choose
\(y\ne x\), \(y\notin\FV(N)\).  Then
\[
\begin{aligned}
 T_{\mathcal B}((\lambda y.R)[N/x])
 &=\mathcal B_y(T_{\mathcal B}(R[N/x]))\\
 &=\mathcal B_y(T_{\mathcal B}(R)[T_{\mathcal B}(N)/x])\\
 &=T_{\mathcal B}(\lambda y.R)[T_{\mathcal B}(N)/x],
\end{aligned}
\]
using the induction hypothesis, variable preservation, and host commutation.
The extensional version for \(T_{\mathcal A}\) uses the same induction with
extensional commutation and abstraction congruence.  The extensional proof
below establishes that congruence directly from the application property before
using any claim about compilation or substitution of lambda terms.

Induction on \(M\) gives \(E(T_{\mathcal D}(M))=_\beta M\): use the expansion
formula in the abstraction case and beta congruence in the application case.
This round trip does not require compiler substitution commutation.
Expansion of a K contraction performs two beta contractions; expansion of an
S contraction performs three.  An application context expands to an application
context, proving the assertion for each weak step.
The root beta simulation is the application property with
\(P=T_{\mathcal D}(M)\), \(Q=T_{\mathcal D}(N)\).  The substitution results
identify its target exactly as stated for each compiler.  Induction on the
application-context grammar propagates the reduction.
\end{proof}
\end{annexitem}

\begin{annexitem}{annex:p07.combinatory.confluence}{thm:p07-combinatory-confluence}{Parallel development of the two contractions}
\begin{proof}
Use the four parallel clauses on the auxiliary sheet.  Induction proves
\(P\Rightarrow_{\mathsf{SK}}P\).  Each weak step is parallel, by identity
premises at its redex and parallel application through its context.  Each
parallel step is a finite weak reduction: reduce its premises by induction,
then perform its root K or S contraction, or retain the application in the
application clause.  Their reflexive transitive closures are thus equal.
Define complete development by
\[
\begin{aligned}
 d(a)&=a\quad(a\text{ an atom}),\\
 d(\mathsf KUV)&=d(U),\\
 d(\mathsf SUVW)&=d(U)d(W)(d(V)d(W)),\\
 d(UV)&=d(U)d(V)\quad\text{if neither preceding root pattern applies}.
\end{aligned}
\]
We prove \(P\Rightarrow_{\mathsf{SK}}Q\) implies
\(Q\Rightarrow_{\mathsf{SK}}d(P)\), by induction on \(P\).
At an atom only identity is possible.  If the root has neither full-arity
pattern, a parallel derivation must use application.  Induction on both children
and parallel application give \(d(P)\), even if the intermediate term has
acquired a new root redex.
For \(P=\mathsf KUV\), a root contraction gives a parallel reduct of \(U\),
which reaches \(d(U)\) by induction.  An application derivation instead gives
\(\mathsf KU'V'\), with parallel reducts of \(U,V\), since the partial
application \(\mathsf KU\) has no root contraction.  Induction and parallel
K reach \(d(U)\).
For \(P=\mathsf SUVW\), a root contraction gives \(U'W'(V'W')\).
Induction and parallel application send both copies of \(W'\) to \(d(W)\),
giving \(d(P)\).  An application derivation gives \(\mathsf SU'V'W'\), because
the partial S applications have no root contraction; induction and parallel S
again give \(d(P)\).  This exhausts all parallel derivations.

Any two parallel reducts have the common parallel reduct \(d(P)\).
Place two finite parallel paths on adjacent sides of a rectangle and fill its
cells using this diamond property.  Induction on rows and columns constructs
the common endpoint; zero-length sides use identity.  The equality of closures
therefore proves weak confluence.  Induction on a conversion zigzag proves its
endpoints joinable: extending forward uses confluence with the existing join;
extending backward gives a path from the new endpoint through the previous one
to that join.  Normal forms have no proper reducts, so two convertible normal
forms are identical.
\end{proof}
\end{annexitem}

\begin{annexitem}{annex:p07.combinatory.weak-boundary}{thm:p07-combinatory-weak-boundary}{A normal-form obstruction beneath a lambda binder}
\begin{proof}
The inner application \((\lambda y.y)x\) contracts to \(x\) beneath the outer
binder.  The host compiler sends \(\lambda y.y\) to I and abstracts
\(\mathsf I x\) as \(\mathsf S(\mathsf K\mathsf I)\mathsf I\).
The term \(\mathsf I=\mathsf S\mathsf K\mathsf K\) is normal: S has only
two arguments and K has none.  In the second output, outer S has two arguments,
K has one, and both I subtrees are normal.  These distinct normal forms are
not weakly convertible by confluence.  Applied to \(Q\), I takes two
contractions to \(Q\).  The second output contracts first to
\((\mathsf K\mathsf I Q)(\mathsf I Q)\).  Contracting K and the right I gives
\(\mathsf I Q\), then \(Q\), proving the application assertion.
\end{proof}
\end{annexitem}

\begin{annexitem}{annex:p07.combinatory.extensional}{thm:p07-combinatory-extensional}{Fresh-variable extensionality and inverse translations}
\begin{proof}
Extensional equality is stable under substitution, by induction on its
generation.  S--K equations, equivalence steps, and application congruence are
directly stable.  At an extensionality step, rename its fresh auxiliary variable
throughout its premise derivation away from the finitely many variables of the
substituting terms.  Renaming a derivation is valid by induction on rule instances:
all side conditions are variable-absence conditions.  This renaming fixes the
conclusion because the auxiliary variable is absent there.  Substitution into
the premise now preserves its freshness, so the rule gives the substituted
conclusion.

If \(P=_{\mathrm e}Q\), choose \(z\) fresh for \(P,Q,x\).
Substitution stability and the application theorem give
\[
 \mathcal D_x(P)z=_{\mathrm e}P[z/x]
   =_{\mathrm e}Q[z/x]=_{\mathrm e}\mathcal D_x(Q)z.
\]
Neither abstractor contains \(z\), so extensionality proves abstraction
congruence, independently of every compilation claim.
The structural substitution induction for \(T_{\mathcal A}\) can therefore
use that congruence at outer abstractors and the preceding theorem's
extensional commutation at a substituted body.  It proves the general
extensional substitution assertion used below.

Check preservation of the generators of beta-eta equality.  A root beta
equation follows from the application theorem and compiler substitution
(literal for the host, extensional for the atom-by-atom algorithm).
Application congruence is part of the equality theory; lambda congruence uses
the abstraction congruence just proved.  For eta, with \(x\notin\FV(M)\),
apply \(\mathcal D_x(T_{\mathcal D}(M)x)\) and \(T_{\mathcal D}(M)\) to a
fresh \(z\).  Both equal \(T_{\mathcal D}(M)z\) by the application theorem;
extensionality proves their equality.  Alpha invariance follows from atomic
renaming.  Equivalence rules now propagate the generator checks to every
conversion derivation.

Expansion sends the S--K equations to beta equality and preserves application
and equivalence rules.  At extensionality, an expanded premise is
\(E(P)z=_{\beta\eta}E(Q)z\) with fresh \(z\).  Abstraction congruence in
lambda calculus gives equality of \(\lambda z.E(P)z\) and
\(\lambda z.E(Q)z\); eta contracts them to \(E(P)\) and \(E(Q)\).
Thus expansion sends extensional equality to beta-eta equality.

For the combinatory round trip, variables are unchanged and application uses
induction.  Apply \(T_{\mathcal D}(E(\mathsf K))\) to fresh \(u,v\).
Successive applications give \(u\), just as \(\mathsf Kuv\) does.
Extensionality first for \(v\), then \(u\), proves equality with K.
On fresh \(u,v,w\), the analogous calculation for S gives \(uw(vw)\);
three extensionality steps prove equality with S.
Now \(E(P)=_{\beta\eta}E(Q)\) implies equality of their compilations and,
by these round trips, \(P=_{\mathrm e}Q\).  Conversely expansion is sound.
Likewise equality of compiled lambda terms implies beta-eta equality of their
expansions and then of the original terms by the beta-correct round trip.
Finally both abstractors applied to a fresh variable reduce to the same
substituted body, so one extensionality step equates them.
\end{proof}
\end{annexitem}

\begin{annexitem}{annex:p07.combinatory.eta-boundary}{thm:p07-combinatory-eta-boundary}{The extra optimization case and its exact limit}
\begin{proof}
Induct on the input tree.  Identity, absent-variable, and application cases
repeat the host proofs.  In the extra case \(P=Ux\), with
\(x\notin\FV(U)\), the output \(U\) has precisely the required free variables
and \(UQ=(Ux)[Q/x]\), with zero reduction steps.  Also
\(E(U)=_{\beta\eta}\lambda x.E(U)x\) by freshness.  In the recursive
application case the expansion proof uses the induction hypotheses as
beta-eta equations, retaining valid beta contraction and congruence.  Structural
induction on lambda terms gives the stated compiled round trip.
The distinct terms \(f\) and \(\lambda x.fx\) are beta-normal when \(x\ne f\).
Beta confluence rules out beta conversion between them.  They are eta-equal,
which identifies the additional equality used by the optimization.
Without freshness, \(xx\) would be sent to \(x\), retaining the eliminated
variable.  On a distinct variable argument \(q\), this gives \(xq\) instead
of the required \(qq\).  These are distinct weak normal forms, so even weak
conversion cannot equate the failed application results.
\end{proof}
\end{annexitem}

\begin{annexitem}{annex:p07.combinatory.completeness}{thm:p07-combinatory-completeness}{Coefficients and all assignment cases}
\begin{proof}
If the structure is combinatorially complete, represent the polynomial
\((xz)(yz)\) in three variables by \(s\), and the polynomial \(x\) in two
indicated variables \(x,y\) by \(k\).  Their representing equations are
exactly the required S and K equations for every assignment.
Conversely assume those equations.  Extend either algorithm to coefficient
atoms by using \(ka\) and use \(i=skk\) in the identity case.
For the assignment identity, induct on the polynomial.  At \(x\), \(ib=b\).
At another variable \(y\), \((k\rho(y))b=\rho(y)\).  At a coefficient \(a\),
\((ka)b=a\).  In the host's whole-subtree absent case, the value is
\((k\llbracket p\rrbracket_\rho)b=\llbracket p\rrbracket_\rho\);
an \(x\)-free tree is unaffected by updating \(\rho(x)\), by induction on
that tree.  For \(uv\), the S equation gives
\[
 (s\llbracket\mathcal D_x(u)\rrbracket_\rho
    \llbracket\mathcal D_x(v)\rrbracket_\rho)b
 = (\llbracket\mathcal D_x(u)\rrbracket_\rho b)
   (\llbracket\mathcal D_x(v)\rrbracket_\rho b)
 =\llbracket uv\rrbracket_{\rho[x\mapsto b]},
\]
the last step being the two induction hypotheses.
For variables \(x_1,\ldots,x_n\), evaluate the closed polynomial formed by
abstracting them in reverse order.  Successive applications of the assignment
identity give its representing equation for every tuple of elements.
For \(n=0\), direct evaluation suffices.  This argument permits all coefficients
throughout and does not require an element argument to have a variable name.

In an extensional combinatory algebra, the S--K equations, their equivalence
consequences, and application congruence hold under every assignment.  For a
fresh-variable extensionality instance, fix every variable other than its
fresh \(z\) and vary \(\rho(z)\) over all elements of the algebra.
The premise says that two fixed elements agree on every argument.
Extensionality gives their equality.  Induction on the equality derivation
proves soundness for all assignments; the syntactic correspondence supplies
the beta-eta consequence for compiled terms.
\end{proof}
\end{annexitem}

\begin{annexitem}{annex:p07.combinatory.size}{thm:p07-combinatory-size}{Tree recurrences and effective compilation}
\begin{proof}
Identity has five nodes; an absent-variable output has \(s(P)+2\leq4s(P)+1\)
nodes; other-atom outputs have three.  In an application case the output size is
\[
 3+s(\mathcal D_x(U))+s(\mathcal D_x(V))
 \leq3+(4s(U)+1)+(4s(V)+1)=4s(UV)+1.
\]
The eta case has size \(s(U)<s(Ux)\).  Since \(s(P)\geq1\), the second
single-abstraction inequality follows.  If every leaf is \(x\), neither
\(\mathcal A\) nor \(\mathcal B\) uses an absent case: equality holds at
each leaf and is preserved by the displayed application recurrence.
For compilation induct on the lambda term.  Variables have size one and no
binders.  For \(UV\), with \(r=b(U)\), \(t=b(V)\),
\[
 1+s(T_{\mathcal D}(U))+s(T_{\mathcal D}(V))
 \leq1+5^r s(U)+5^t s(V)
 \leq5^{r+t}(1+s(U)+s(V)).
\]
For \(\lambda x.U\), the single-abstraction estimate gives
\(s(T_{\mathcal D}(\lambda x.U))\leq5s(T_{\mathcal D}(U))
\leq5^{b(U)+1}s(U)\leq5^{b(\lambda x.U)}s(\lambda x.U)\).
Expanded K has three nodes and expanded S has ten; expansion fixes variable
leaves and application nodes.  Induction proves \(s(E(P))\leq10s(P)\).

For one abstraction, mark occurrence of \(x\) in every subtree in one
bottom-up pass.  The output pass uses these Boolean marks for constant-time
branch decisions; the eta case additionally checks whether the right child is
\(x\).  Each visited node causes bounded constructor work beyond recursive
calls.  An absent subtree copied into the output is copied once and the
recursion stops there; such subtrees are disjoint.  Charging to input visits
and constructed output nodes proves linear time with unit-cost node operations
and variable comparisons.  Repeated binders operate on enlarged intermediate
trees; this calculation asserts no polynomial bound for the entire compilation.
\end{proof}
\end{annexitem}
